\section{Group fields}\label{sec:groupFields}

A group field is a vector
\begin{align*}
    \hypercoreat{\catvariable} \wcols G \rightarrow \mathbb{C}
\end{align*}
where $\catvariable$ enumerated the $\cardof{G}=\catdim$ elements of the group.
The $\catindex$-th element of the group is then mapped to $\hypercoreat{\indexedcatvariable}$.

\begin{example}{Unfolding is not an isomorphism}
    Let us note that the unfolding of $[\prod_{\catenumeratorin}\catdimof{\catenumerator}]$ to $\facstates$ is not a group isomorphism.
    The group operation in $\prod_{\catenumeratorin}\catdimof{\catenumerator}$ is the integer summation
    \begin{align*}
        \catindex \circ \seccatindex =
        \catindex + \seccatindex \modspace \left(\prod_{\catenumeratorin}\catdimof{\catenumerator}\right) \, .
    \end{align*}
    In the unfolding the group operation is
    \begin{align*}
        (\catindexof{0},\ldots,\catindexof{\catorder-1}) \circ (\seccatindexof{0},\ldots,\seccatindexof{\catorder-1})
        = (\catindexof{0}+\seccatindexof{0}\modspace \catdimof{0},\ldots,\catindexof{\catorder-1}+\seccatindexof{\catorder-1}\modspace \catdimof{\catorder-1}) \, .
    \end{align*}
    Intuitively, the difference of these groups is that in the unfolded group summation carry bits are set to $0$.

\end{example}

The group field perspective on tensors further applied in quantum gravity.
Here symmetries are modeled based on the group structure \cite{chirco_group_2018,chirco_group_2020}.

\subsection{Coset encoding}

We here restrict to finite abelian groups.
A subgroup is a subset $H\subset G$, which is itself a group.
A coset to $H$ and an element $g \in G$ is the set
\begin{align*}
    \{g\circ h \wcols h\in H\} \, .
\end{align*}

% Dual group
The dual group $\hat{G}$ is the collection of all characters of irreducible representations
\begin{align*}
    \chi_{k} \defcols G \rightarrow \mathbb{C} \, ,
\end{align*}
with the group operation $\chi_{k_0}\circ \chi_{k_1} (g) = \chi_{k_0}(g)\cdot\chi_{k_1}(g)$.
The annihilator of a subgroup $H\subset G$ is the subgroup
\begin{align*}
    H^{\perp}\coloneqq
    \{\chi_{k} \wcols \chi_{k}\in\hat{G}\ncond \uniquantwrtof{h\in H}{\chi_{k}(h)=1} \} \, .
\end{align*}
The annihilator of the annihilator is
\begin{align*}
    \big(H^{\perp}\big)^{\perp} = H \, .
\end{align*}

Since finite Abelian group are isomorphic to cartesian products of cyclic groups, we assume $G=\bigtimes_{\catenumeratorin}[\catdimof{\catenumerator}]$.
The characters are enumerated by $\thirdcatindexof{[\catorder]}\in\bigtimes_{\catenumeratorin}[\catdimof{\catenumerator}]$ and defined as the maps
\begin{align*}
    \chi_{\thirdcatindexof{[\catorder]}}(\shortcatindices)
    = \expof{\twopii \sum_{\catenumeratorin} \frac{\thirdcatindexof{\catenumerator}\cdot\catindexof{\catenumerator}}{\catdimof{\catenumerator}}} \, .
\end{align*}
The annihilator of a subgroup generated by $h^j_{[\catorder]}$ for $j\in[n]$ is thus defined by
\begin{align*}
    \{\chi_{\thirdcatindexof{[\catorder]}} \wcols
    \uniquantwrtof{j\in[n]}{
        \sum_{\catenumeratorin} \frac{\thirdcatindexof{\catenumerator}\cdot h^j_{\catenumerator}}{\catdimof{\catenumerator}} = 0 \modspace 1
    }\} \, .
\end{align*}


We build for each $h^j_{[\catorder]}$ an edge
\begin{align*}
    \edgeof{j}
    =\{\catenumerator \wcols h^j_{\catenumerator}\neq 0\}
\end{align*}
and a tensor
\begin{align*}
    \hypercoreofat{h^j_{[\catorder]}}{\indexedcatvariableof{\edgeof{j}}}
    =\begin{cases}
         1 & \ifspace \sum_{\catenumerator\in\edgeof{j}} \frac{\thirdcatindexof{\catenumerator}\cdot h^j_{\catenumerator}}{\catdimof{\catenumerator}} = 0 \modspace 1  \\
         0 & \elsetext
    \end{cases} \, .
\end{align*}
Then we have for the subset encoding of $H^{\perp}\subset\hat{G}$ (where $\hat{G}$ is enumerated by $\shortcatindicesin$) that
\begin{align*}
    \onehotmapofat{H^{\perp}}{\shortcatvariables}
    = \contractionof{
        \{\hypercoreofat{h^j_{[\catorder]}}{\indexedcatvariableof{\edgeof{j}}} \wcols j\in[n]\}
    }{\shortcatvariables} \, .
\end{align*}

\begin{example}\label{exa:cosetsCyclicTwoProduct} % WH case
    If $\catdimof{\catenumerator}=2$ we can simplify the support condition for $\hypercoreofat{h^j_{[\catorder]}}{\indexedcatvariableof{\edgeof{j}}}$, since any nonvanishing $h^j_{\catenumerator}$ is $1$.
    The condition becomes
    \begin{align*}
        \sum_{\catenumerator\in\edgeof{j}} \thirdcatindexof{\catenumerator} = 0 \modspace 2
    \end{align*}
    Notice, that this is captured by the generalized exactly or connective and thus
    \begin{align*}
        \hypercoreofat{h^j_{[\catorder]}}{\catvariableof{\edgeof{j}}}
        = \bigoplus [\catvariableof{\edgeof{j}}] \, .
    \end{align*}
    Using the associativity of the exactly or connective, one can decompose the tensor in an arbitrary tree format with hidden ranks $2$.
\end{example}

% Representing different
Having represented a group itself based on its annihilator, we have represented the coset of $H$ to the neutral group element $0$.
To represent an arbitrary coset to an element $\thirdcatindexof{[\catorder]}$ we apply a shift operator
\begin{align*}
    \bigotimes_{\catenumeratorin} \hypercoreofat{\catenumerator}{\catvariableof{\catenumerator},\tildecatvariableof{\catenumerator}} \, ,
\end{align*}
where each local shift operator is a basis encoding of a shift function defined as
\begin{align*}
    \hypercoreofat{\catenumerator}{\indexedcatvariableof{\catenumerator},\tildecatvariableof{\catenumerator}=\tildecatindexof{\catenumerator}}
    = \begin{cases}
          1 & \ifspace \catindexof{\catenumerator} + \thirdcatindexof{\catenumerator} = \tildecatindexof{\catenumerator} \modspace \catdimof{\catenumerator}\\
          0 & \elsetext
    \end{cases} \, .
\end{align*}